% Part: proof-theory
% Chapter: normalization
% Section: introduction

\documentclass[../../../include/open-logic-section]{subfiles}

\begin{document}

\olfileid{pt}{nor}{seg}

\olsection{Segmen lan Cut}

Aturan !!^a{introduction} kang disusul aturan !!a{elimination}
iku liku-liku ing !!a{proof}; normalisasi !!a{proof} tegesé
ngilangi liku-liku kuwi. Nanging, \Elim{\lor} lan
\Elim{\lexists} ndadèkaké definisi ``liku-liku kang arep
diilangi'' rada ruwet. Masalahé mangkéné. Amarga aturan iki,
ing liku-liku ora mesthi aturan !!{elimination} \emph{langsung}
nyusul aturan !!{introduction}. Ing antarané bisa ana
\Elim{\lor} utawa \Elim{\lexists}, upamané,
\[
\AxiomC{}
\DeduceC{$\lexists[x][!C(x)]$}
\AxiomC{$\Discharge{!A}{x}$}
\DeduceC{$!B$}
\DischargeRule{\Intro{\lif}}{x}
\UnaryInfC{$!A \lif !B$}
\DischargeRule{\Elim{\lexists}}{y}
\BinaryInfC{$!A \lif !B$}
\AxiomC{}
\DeduceC{$!A$}
\RightLabel{\Elim{\lif}}
\BinaryInfC{$!B$}
\DisplayProof
\]
Malah, cacah pira waé aturan \Elim{\lor} utawa
\Elim{\lexists} bisa misahaké \Intro{\lif} saka \Elim{\lif},
upamané,
\[
\AxiomC{}
\DeduceC{$\lexists[x][!C(x)]$}
\AxiomC{}
\DeduceC{$!D \lor !E$}
\AxiomC{$\Discharge{!A}{x}$}
\DeduceC{$!B$}
\DischargeRule{\Intro{\lif}}{x}
\UnaryInfC{$!A \lif !B$}
\AxiomC{$\Discharge{!A}{x}$}
\DeduceC{$!B$}
\DischargeRule{\Intro{\lif}}{x}
\UnaryInfC{$!A \lif !B$}
\RightLabel{\Elim{\lor}}
\TrinaryInfC{$!A \lif !B$}
\DischargeRule{\Elim{\lexists}}{y}
\BinaryInfC{$!A \lif !B$}
\AxiomC{}
\DeduceC{$!A$}
\RightLabel{\Elim{\lif}}
\BinaryInfC{$!B$}
\DisplayProof
\]
Tuladha iki nuduhaké yèn siji \Elim{\lif} bisa kalebu ing
pirang-pirang liku-liku, amarga inferènsi mau nyusul luwih saka
siji~\Intro{\lif}. Kanggo nyakup kabèh kemungkinan iki, kita
netepaké liku-liku nganggo urutan tartamtu saka kedadéan
!!{formula} ing bukti~\Log{N1i} kang diarani
\emph{segmen}. Ing tuladha iki, urutan mau ngemot kedadéan
~$!A \lif !B$, nalika premis tambahan saka $\Elim{\lor}$ utawa
$\Elim{\lexists}$ nyusul konklusiné. Ana rong segmen kaya iki,
saben-saben ngemot telung kedadéan~$!A \lif !B$.
Definisi resminé mangkéné.

\begin{defn}
\emph{Segmen} ing \Log{N1i} !!{proof}~$\delta$ iku urutan
kedadéan $!A_1$, \dots, $!A_n$ saka !!{formula}~$!A$ kang padha
ing $\delta$, kanthi syarat
\begin{enumerate}
  \item $!A_1$ dudu konklusi saka \Elim{\lor} utawa \Elim{\lexists};
  \item $!A_n$ dudu premis tambahan saka \Elim{\lor} utawa
  \Elim{\lexists};
  \item Kanggo $i = 1$, \dots,~$n-1$, $!A_i$ iku premis tambahan
  saka \Elim{\lor} utawa \Elim{\lexists}, lan $!A_{i+1}$
  iku konklusi inferènsi mau.
\end{enumerate}
\emph{Dawa} segmen mau yaiku~$n$.
\end{defn}

Ora saben segmen dadi sasaran kanggo diilangi (yaiku liku-liku).
Segmen kang dadi sasaran diarani \emph{cut}.

\begin{defn}
\emph{Cut} iku segmen kang $!A_n$ dadi premis utama inferènsi
!!a{elimination}, lan salah siji saka loro kahanan iki kelakon:
$n=1$ lan $!A_1 = !A_n$ dadi konklusi inferènsi
!!a{introduction} utawa~\FalseInt; utawa~$n>1$.
\end{defn}

Elinga yèn segmen cut ora kudu diwiwiti saka konklusi aturan
!!a{introduction}. Cukup yèn pungkasané dadi premis utama
aturan !!a{elimination}. Nanging, supaya segmen dawa~$1$
kacathet minangka cut, segmen mau kudu dadi konklusi aturan
!!{introduction} utawa~\FalseInt.

\begin{defn}
\emph{Pangkat} cut yaiku~$\depth{!A}$. \emph{Pangkat
cut}~$\cutrank{\delta}$ saka !!a{proof}~$\delta$ iku pangkat
paling gedhé ing antarané cut-cut ing~$\delta$. Cut diarani
\emph{maksimal} ing~$\delta$ yèn pangkaté
padha karo~$\cutrank{\delta}$, yaiku pangkat maksimal.
\emph{Dawa cut}~$\delta$ iku gunggung dawa kabèh cut maksimalé.
\end{defn}

\end{document}
